\documentclass[11pt]{article}
\usepackage[utf8]{inputenc}
\usepackage{amsmath}
\usepackage{amssymb}
\usepackage{amsfonts}
\usepackage[margin=0.6in, top=0.5in]{geometry}
\usepackage{enumitem}
\usepackage{chemfig}
\usepackage{array}
\usepackage{tabularx}

% Custom command for the communication/marking header box
\newcommand{\testheader}[1]{
\noindent
\begin{tabularx}{\textwidth}{|X|m{4.5cm}|}
\hline
\textbf{\underline{Communication:}} & \textbf{\underline{Overall:}} \hfill #1 of 5 \\
\textbf{Language:} easy to understand, answers question, follows instructions, concise, precise & \\
\textbf{Chemistry:} proper vocab, notation, spacing, logical & Basic / ok / good / clearly written \\
\textbf{Math:} organized, units, sig digs, math notation, concluding statement, spacing & \space\space L1 \space\space L2 \space\space L3 \space\space\space\space\space L4 \\
\hline
\end{tabularx}
\vspace{0.4cm}
}

\begin{document}

% --- PAGE 1 ---
\testheader{1}

\noindent
SCH4U \hfill Name: \underline{\hspace{5cm}} \\
\indent \hfill Tuesday, April 2, 2025

\begin{center}
    \Large \textbf{Unit 2: Rates \& Energy Unit Test}
\end{center}

\noindent
\begin{tabular}{|l r|l r|l r|}
\hline
Total: & /37 & Comm: & /10 & LL: & /5 \\
\hline
\end{tabular}

\vspace{0.4cm}

\noindent
\underline{\textbf{Specific Heat Capacities of Water:}} \space\space\space\space $c_{(s)} = 2.01\text{ J/g}^\circ\text{C}$, $c_{(l)} = 4.184\text{ J/g}^\circ\text{C}$, $c_{(g)} = 2.01\text{ J/g}^\circ\text{C}$ \\
\underline{\textbf{Molar Enthalpies for Methanol:}} \space\space\space\space\space $\Delta H_{fus} = 3.20\text{ kJ/mol}$; $\Delta H_{vap} = 35.2\text{ kJ/mol}$

\vspace{0.5cm}

\begin{enumerate}[leftmargin=*, label=\arabic*.]
    \item (2 marks) A chemist boiled 0.300 L of water (at $100^\circ\text{C}$) \& had just transferred it into an insulated container. Unfortunately, while looking down at the water, the chemist's 40.0 g gold necklace broke \& fell into the water. Assuming the initial temperature of the gold necklace was at $37.0^\circ\text{C}$, set up the expression (but do not solve or isolate the variable) to find the final temperature of the necklace \& the water. ($c_{water} = 4.184\text{ J/g}^\circ\text{C}$, $c_{gold} = 0.129\text{ J/g}^\circ\text{C}$)
    \vspace{4.5cm}

    \item (2 marks) A student decides to perform a calorimetry experiment in order to calculate the specific heat of a metal. The student heats the metal and uses a glass stirring rod to stir the contents of the experiment before using the thermometer to determine the final temperature. How would this action affect the numerical results of the experiment? Briefly explain.
    \vspace{3.5cm}

    \item (3 marks) What are the three criteria for an effective collision (point form)?
\end{enumerate}

\vfill

\newpage

% --- PAGE 2 ---
\testheader{2}

\begin{enumerate}[leftmargin=*, label=\arabic*., start=4]
    \item Suppose a chemical reaction occurs as follows: $\text{A (s)} + \text{B (g)} + \text{C (g)} \rightarrow \text{D (g)} + \text{E (g)}$. Answer the following.
    \begin{enumerate}[label=\alph*.)]
        \item (1 mark) Do we need to know what the products are in determining the reaction rate? Explain.
        \vspace{2.5cm}
        \item (1 mark) Looking at the overall reaction, do you know if the reaction consists of a single elementary step or will it require a multi-step reaction mechanism? Explain. DO NOT REFER TO PART (E) WHEN ANSWERING THIS QUESTION.
        \vspace{2.5cm}
        \item (3 marks) From the following data, what is the rate law for the reaction? \\
        
        \begin{center}
        \begin{tabular}{|c|c|c|c|c|}
        \hline
        Trial & [A] initial (mol/L) & [B] initial (mol/L) & [C] initial (mol/L) & Initial rate (mol/L$\cdot$s) \\ \hline
        1 & 0.10 & 0.10 & 0.30 & 0.00126 \\ \hline
        2 & 0.10 & 0.10 & 0.60 & 0.00504 \\ \hline
        3 & 0.10 & 0.20 & 0.30 & 0.00252 \\ \hline
        4 & 0.20 & 0.10 & 0.30 & 0.00126 \\ \hline
        \end{tabular}
        \end{center}
        \vspace{4cm}
        \item (2 marks) What is the value of the rate law constant for the reaction?
    \end{enumerate}
\end{enumerate}

\vfill

\newpage

% --- PAGE 3 ---
\testheader{3}

\begin{enumerate}[leftmargin=*, label=\arabic*., start=5]
    \item Copper (II) hydroxide, $\text{Cu(OH)}_2$, is a solid at SATP conditions.
    \begin{enumerate}[label=\alph*.)]
        \item (1 mark) Write out the formation equation for copper (II) hydroxide.
        \vspace{1.5cm}
        \item (2 marks) Using Collision Theory, explain what happen to the overall rate if the temperature of the system was increased.
        \vspace{3.5cm}
        \item (2 marks) Draw the diagram for the reaction under two different temperature conditions (the reaction is exothermic). Indicate the activation energy.
        \vspace{0.5cm}
        
        \begin{tabular}{cc}
        \begin{minipage}{0.45\textwidth}
        \begin{picture}(160,130)
        \put(20,20){\vector(1,0){130}}
        \put(20,20){\vector(0,1){110}}
        \put(5,115){\rotatebox{90}{Energy}}
        \put(55,5){\# of molecules}
        \end{picture}
        \end{minipage}
        &
        \begin{minipage}{0.45\textwidth}
        \begin{picture}(160,130)
        \put(20,20){\vector(1,0){130}}
        \put(20,20){\vector(0,1){110}}
        \put(5,115){\rotatebox{90}{Energy}}
        \put(50,5){reaction coordinate}
        \end{picture}
        \end{minipage}
        \end{tabular}
    \end{enumerate}

    \vspace{1cm}

    \item The forward $E_A$ for an elementary reaction is $42\text{ kJ}$ \& the reverse $E_A$ is $32\text{ kJ}$
    \begin{enumerate}[label=\alph*.)]
        \item (2 marks) Is the reaction exothermic or endothermic? Explain.
        \vspace{2.5cm}
        \item (1 mark) Calculate $\Delta H$.
    \end{enumerate}
\end{enumerate}

\vfill

\newpage

% --- PAGE 4 ---
\testheader{4}

\begin{enumerate}[leftmargin=*, label=\arabic*., start=7]
    \item Methanol, $\text{CH}_3\text{OH}$, has a melting point temperature of $-98^\circ\text{C}$ \& a boiling point temperature of $65.0^\circ\text{C}$. It has a molar mass of $32.04\text{ g/mol}$. Refer to the reference material for additional relevant data.
    \begin{enumerate}[label=\alph*.)]
        \item (2 marks) Draw a Temperature/Enthalpy graph when $10.0\text{ g}$ of methanol vapour at $65.0^\circ\text{C}$ cools \& forms liquid methanol at $-90.0^\circ\text{C}$
        \vspace{5cm}
        \item (4 marks) Calculate the enthalpy change. Show your work.
        \vspace{4.5cm}
    \end{enumerate}

    \item (4 marks) Calculate the enthalpy for this reaction: $2\text{C}_{(s)} + \text{H}_{2(g)} \rightarrow \text{C}_2\text{H}_{2(g)} \text{ } \Delta H^\circ=?$
    \vspace{0.3cm}
    
    \begin{enumerate}[label=(\roman*)]
        \item $\text{C}_2\text{H}_{2(g)} + \frac{5}{2}\text{O}_{2(g)} \rightarrow 2\text{CO}_{2(g)} + \text{H}_2\text{O}_{(l)} \hfill \Delta H^\circ = -1299.5\text{ kJ}$
        \item $\text{C}_{(s)} + \text{O}_{2(g)} \rightarrow \text{CO}_{2(g)} \hfill \Delta H^\circ = -393.5\text{ kJ}$
        \item $\text{H}_{2(g)} + \frac{1}{2}\text{O}_{2(g)} \rightarrow \text{H}_2\text{O}_{(l)} \hfill \Delta H^\circ = -285.8\text{ kJ}$
    \end{enumerate}
\end{enumerate}

\vfill

\newpage

% --- PAGE 5 ---
\testheader{5}

\begin{enumerate}[leftmargin=*, label=\arabic*., start=9]
    \item (4 marks) Ammonia, $\text{NH}_3$, reacts with oxygen:
    \begin{center}
        $4\text{NH}_{3(g)} + 5\text{O}_{2(g)} \rightarrow 4\text{NO}_{(g)} + 6\text{H}_2\text{O}_{(l)}$
    \end{center}
    Calculate the enthalpy of reaction if $0.0064\text{ g}$ of ammonia reacts completely.
    \begin{flushright}
        \begin{tabular}{|c|c|c|}
        \hline
        \multicolumn{2}{|c|}{} & $\Delta H_f$ \\ \hline
        ammonia (aq) & $\text{NH}_{3(aq)}$ & $-80.0\text{ kJ/mol}$ \\ \hline
        ammonia (g) & $\text{NH}_{3(g)}$ & $-45.9\text{ kJ/mol}$ \\ \hline
        nitrogen dioxide & $\text{NO}_{2(g)}$ & $33.2\text{ kJ/mol}$ \\ \hline
        nitrogen monoxide & $\text{NO}_{(g)}$ & $90.2\text{ kJ/mol}$ \\ \hline
        water & $\text{H}_2\text{O}_{(g)}$ & $-241.8\text{ kJ/mol}$ \\ \hline
        water & $\text{H}_2\text{O}_{(l)}$ & $-285.8\text{ kJ/mol}$ \\ \hline
        \end{tabular}
    \end{flushright}
    \vspace{7cm}
    \begin{enumerate}[label=\alph*.)]
        \item (1 mark) List all reaction intermediates.
    \end{enumerate}
\end{enumerate}

\vfill

\end{document}